\begin{abcliste}
\abc The electrons get the kinetic energy $\frac{p^2}{2m_e} = eU$, so $p = \sqrt{2 m_e e U}$ and
\begin{align}
 \lambda &= \frac{h}{p} = \laF \\
   &= \frac{\nch}{\sqrt{2\times\ncme\times\nce\times\U}} \\
   &= \la \approx \result{\laP{-9}{2}}
\end{align}
(The formula $\lambda = hc/E$ holds for photons only.)

\abc
\begin{align}
 \varphi &= \phF = \arcsin\frac{\la}{\d} = \ph \approx \result{\phP{0}{2}}
\end{align}
At \UO, Davisson and Germer found the maximum at \SI{50}{\degree}: the first proof that electrons are waves.

\abc A higher voltage gives faster electrons with a shorter wavelength, so $\sin\varphi = \lambda/d$ gets smaller: the maximum moves $\result{\text{to smaller angles}}$.
\end{abcliste}
